\section{JASPによるベイズ推定の実践}
\label{s28_jasp}

\subsection{授業内容}
\subsubsection{科目の中でのこのコマの位置づけ}
\begin{tabular}[t]{cccccc}
	記述統計[3] & 確率[14] & モデル[15] & 推測・推定[9] & 検定[6] & 実践[6]
\end{tabular}
\subsubsection{概要}
ベイズ統計を用いた考え方は，従来の推定やモデル比較と異なる観点からのアプローチができるようになる。もっとも原理的に優れたものであっても，実践できないのであれば絵に描いた餅にすぎない。
JASPの登場により，GUIで比較的簡単にベイズ法による分析，実践を行うことが可能になった。ここではJASPの導入と，JASPによる回帰分析，帰無仮説検定を行い，従来法とベイズ統計とでどのような違いがあるのかを演習的に学ぶ。
なお，JASPのバージョンは本シラバス執筆時点\footnote{2021年12月11日}でver.0.16であり，未だ発展途上で急速に開発が進んでいる段階であるため，表示や操作の方法が今後変更になる可能性が高いことに留意する。
\subsubsection{コマ主題細目(箇条書き)}

\begin{description}
	\item[ベイズ推定の方法] ベイズ推定を実際に計算するためには，解析的な方法やグリッドサーチ，そしてMCMCによる推定方法が必要になる。とくにマルコフチェイン・モンテカルロ法は，困難であったベイズ推測を可能にした技術的ブレイクスルーであった。基本的なアイデアは，確率分布の形状を求める際に乱数を使って近似するということである。単純な乱数を用いた例から，これによる確率分布の近似と積分計算の簡略化ができることを理解し，事後分布の相対的な形状近似ができるようになったことを学ぶ。その上で，便利なアプリケーションであるJASPの紹介を行う。
	      \begin{flushright}
		      $\to$ \citet{Maeda2017}の7章，\citet{Baba201907}P.62--76
	      \end{flushright}
	\item[JASPの導入] JASPのインストールと起動，その基本的な特徴の解説など導入する演習を行う。各自のPCにJASP環境を導入し，サンプルコードを開いて記述統計量やプロットを描く。現段階\footnote{2021.12.01現在}ではJASPを用いたテキストや解説書は少ないが，標準的なアプリケーションの導入で対応できるため，初学者にもその障壁は低いと考えられる。
	\item[JASPによる回帰分析] JASPの特徴の一つは，従来法とベイズ法が同じメニューから並列的に利用できることである。ここでは回帰分析を例に，最尤法とベイズ法で結果の表示の違いや解釈の方法について見ていく。
	\item[JASPによる平均の比較] 一要因二水準，一要因三水準の平均の比較について，従来法とベイズ法の両方をJASPで実践する。帰無仮説検定における$p$値の解釈と，ベイズ法におけるベイズファクターでの評価それぞれについて，演習的に学ぶ。

\end{description}
\subsubsection{キーワード}
\begin{itemize}
	\item MCMC法
	\item JASP
	\item BF
\end{itemize}
\subsection{授業情報}
\paragraph{コマの展開方法}
演集
\paragraph{標準シラバスにおける位置づけ}
科目番号5；心理学統計法；(1)心理学で用いられる統計手法；J より高度な記述や推定を目指して
\subsubsection{予習・復習課題}
\paragraph{予習}本講では新しい統計ソフト，JASPが導入される。授業開始前にJASPを自らのPCにインストールし，起動するかどうかを確認しておこう。
\paragraph{復習}これまでのデータ分析はRで行ってきたが，同じデータをJASPで行えばどうなるのか，進んで実行してもらいたい。またJASPにはサンプルコードやOSFとの連携機能も含まれているので，これら既存のデータを使ってJASPがどのような振る舞いをするのか，さまざまに試してみると良いだろう。
